\scanpage{12}
\mspage{11}

\deleted{\unsure{Considérons}}
On peut imposer des conditions à $A,B$, comme d'être connexes, voire, contractiles, voire $=$, homéom. à un espace donné, et itou pour les lieux, \inserted{de façon à travailler sur un sous-ens $\mathcal L_P$, $P$ propriété \unsure{intrinsèque}}.

Pour que les conditions plus haut sur l'ens \uline{ordonné} $\mathcal L$ soient encore vraies pour $\mathcal L_P$, il est bon de \deleted{\unsure{outre}} \deleted{de} supposer \deleted{que} (pour $P$ \unsure{choisir|choisie} p.ex. d'être contractile, ou homéomorphe à tel espace fixé) que les espaces $S_A,S_B$ (bords des vois. tub. de $A,B$ dans $X$) aient cette propriété; et \deleted{\unsure{itou}} on fera la restriction \inserted{correspondante} sur les lieux internes $C$: \unsure{les} bords des vois. tub. de $C$ dans $X_{A,C}$ et $X_{B,C}$ ont la propriété $P$.

\uline{Exemples}

1) $P$: être \deleted{homéo} réduit à un pt. On trouve la théorie des segments top. \unsure{ord.|ind.}

\rotatebox{90}{\reflectbox{\unsure{EXTRA STRONG}}}
\rotatebox{90}{\reflectbox{\unsure{EXTRA STRONG}}}
\mspage{11}
